% 图 2.3 —— 手工重建（y = f(x) 的抛物线 + 刻度轴）
%
% 做法：看图定「是什么」，算法在圈定区域内量「是多少」。
%
%   抛物线   不是逼近出来的 —— **认出它就是 y = 3x² + 2**：
%             顶点实测 (356,568) 即数据点 (0,2)；过 (243,229)/(469,229) 即 (±1,5)；
%             沿曲线逐点验证 |x|≤0.8 处到墨迹距离全是 0.00 ✓
%             所以用 plot[domain] 解析画，不再拿描摹点做样条（那样会有波纹）。
%   坐标     原点 (356,794)，比例 **113 px 每单位**（x 刻度实测 124/240/465/577）
%   刻度     两轴各画短线；数字位置取实测框中心
%   空心圆点 (±1,5) 与 (0,2)：外径实测 18、描边 4 → 中线半径 7
%   公式     `y = f(x)` 加**白底**，免得被曲线压住（用户指出）
\documentclass[12pt]{standalone}
\usepackage{fontspec}
\setsansfont{Arial}
\newfontfamily\mfallback{DejaVu Sans}
\usepackage{tikz}

\def\ox{356}       % 原点
\def\oy{794}
\def\U{113}        % 每单位像素
\def\wAx{4.0}      % 轴线宽（实测）
\def\wCur{5.5}     % 曲线宽（实测，比轴粗）

% 抛物线 y = 3x^2 + 2：图像坐标 = (\ox + \U*t, \oy - \U*(3t^2+2))
% 两臂在 t = ±1.24 处出画框顶（实测端点 y≈45）
\begin{document}
\begin{tikzpicture}[x=1pt, y=-1pt]

  % ── 两轴 ──
  \draw[line width=\wAx] (14,794) -- (700,794);     % x 轴
  \draw[line width=\wAx] (356,10) -- (356,890);     % y 轴

  % ── 刻度（短线，长度实测 ≈16）──
  \foreach \t in {-2,-1,1,2}{ \draw[line width=\wAx] ({\ox+\U*\t},786) -- ({\ox+\U*\t},802); }
  \foreach \t in {1,2,3,4,5,6}{ \draw[line width=\wAx] (348,{\oy-\U*\t}) -- (364,{\oy-\U*\t}); }

  % ── 抛物线（解析）──
  \draw[line width=\wCur] plot[domain=-1.24:1.24,samples=161]
       ({\ox+\U*\x},{\oy-\U*(3*\x*\x+2)});

  % ── 三个实心圆点 ──
  % ⚠ 原来画的是**空心环**（\draw + fill=white + circle(7)）—— 那是错的。
  %   原书 600dpi 上逐行量墨段：经过标记中心那一行**只有一段墨**（宽 18 单位），
  %   圆环的话该行必有两段（左右弧）。三个点实测外径 18.3 / 17.2 / 19.4，
  %   故一律 `\fill ... circle (9)`。
  %   附带：空心环的 fill=white 还会把穿过它的抛物线**抠掉一块**，
  %   那正是用户说的「数字 2 挡住曲线、观感不对」的一部分成因。
  \fill (\ox-\U,\oy-5*\U) circle (9);
  \fill (\ox+\U,\oy-5*\U) circle (9);
  \fill (\ox,\oy-2*\U)   circle (9);

  % ── x 轴数字 ──
  \node[anchor=center,inner sep=1.5pt,fill=white,font=\fontsize{28}{35}\selectfont] at (126.0,822.0) {\textsf{\textbf{−2}}};
  \node[anchor=center,inner sep=1.5pt,fill=white,font=\fontsize{28}{35}\selectfont] at (239.5,822.0) {\textsf{\textbf{−1}}};
  \node[anchor=center,inner sep=1.5pt,fill=white,font=\fontsize{28}{35}\selectfont] at (465.0,822.0) {\textsf{\textbf{1}}};
  \node[anchor=center,inner sep=1.5pt,fill=white,font=\fontsize{28}{35}\selectfont] at (579.5,822.0) {\textsf{\textbf{2}}};

  % ── y 轴数字（位置取实测框中心）──
  \node[anchor=center,inner sep=1.5pt,fill=white,font=\fontsize{28}{35}\selectfont] at (337.5, 99.0) {\textsf{\textbf{6}}};
  \node[anchor=center,inner sep=1.5pt,fill=white,font=\fontsize{28}{35}\selectfont] at (337.0,215.5) {\textsf{\textbf{5}}};
  \node[anchor=center,inner sep=1.5pt,fill=white,font=\fontsize{28}{35}\selectfont] at (335.5,332.0) {\textsf{\textbf{4}}};
  \node[anchor=center,inner sep=1.5pt,fill=white,font=\fontsize{28}{35}\selectfont] at (335.5,449.0) {\textsf{\textbf{3}}};
  % ⚠ 用户两次指出「数字 2 挡住曲线」。**白底是用户要的，不能去掉** ——
  %   要去掉的是「白底与曲线/圆点相交」这件事，办法是**把数字挪到不相交的位置**。
  %
  %   节点方框 = 字形 + inner sep 1.5pt ⇒ 半宽约 9.3、半高约 11.5（画布单位）。
  %   两处障碍：① 抛物线 —— 其最低点是顶点画布 y=568（数据 y=2），即曲线只存在于
  %   y≤568；② 顶点圆点 —— 半径 9、圆心 (356,568)，左缘 x=347。
  %   取 (332,578)：方框 x 322.7..341.3、y 566.5..589.5
  %     · 曲线在 y=566.5 处已退到 x≥348.5 > 341.3 ✓
  %     · 圆点左缘 347 > 341.3 ✓
  %   两处都不相交，白底完整，曲线和圆点也不再被切。
  \node[anchor=center,inner sep=1.5pt,fill=white,font=\fontsize{28}{35}\selectfont] at (332.0,578.0) {\textsf{\textbf{2}}};
  \node[anchor=center,inner sep=1.5pt,fill=white,font=\fontsize{28}{35}\selectfont] at (337.0,676.5) {\textsf{\textbf{1}}};

  % ── 公式：**加白底**，免得与曲线/刻度叠在一起看不清 ──
  \node[anchor=center,inner sep=4pt,fill=white,font=\fontsize{26}{32.5}\selectfont] at (533.0,312.0)
       {\textsf{\textbf{\textit{y} = \textit{f}(\textit{x})}}};

  \path[use as bounding box] (0,0) rectangle (710,897);
\end{tikzpicture}
\end{document}
